% Standalone mathematical section; requires amsmath, amssymb, amsthm.
% Original source tasks are not modified by this section.
\section{Exact bridges from the retained Fable polynomial and the S6 cusp}
\label{sec:s6b}

\subsection{Source and scope}
The source corpus is \emph{S6 Proof Audit, Reconstruction \& Overleaf Release};
its audited artifact names and hashes are retained with this notebook.
Its current source root is
\path{s6_consolidated_canonical_2026-08-26}.  All source locators
in this section are relative to that root.  The dated status in
\path{supporting_materials/workbench/research/status.tex},
lines 145--154 and 527--587, says that all six topology bridges are closed and
records a compact complex threefold $X$ with
$X\cong_{\mathrm{diff}}S^6$, $a(X)=1$, and
$\langle c_3(T_X),[X]\rangle=2$.  It explicitly defines ``Verified'' to mean
that a complete proof is written in the workbench.  The older unconfirmed
audit verdict at \texttt{main.tex}, lines 5989 and 6241, is not the current
workbench verdict.  The present bounded audit checks the concrete polynomial,
metric, differential-operator and Fourier calculations below; it does not
independently re-certify every analytic gluing and integral geometric map of
the complete S6 workbench.

For example, the displayed terminal group presentation at status lines
578--580 is trivial by an exact calculation: $xy=1$ gives $y=x^{-1}$;
$x^3=c$ and $y^4=c^{-1}$ give $x^{-4}=x^{-3}$, hence $x=1$, then $y=c=1$.
This verifies the algebra of that presentation.  Whether it is the geometric
fundamental group is a different claim, whose geometric proof is in the
source; this short algebra calculation does not replace that proof.

The polynomial source is
\path{supporting_materials/workbench/comparative_context/prior_fable_context.tex},
lines 8--28.  That file describes a pinned JGPTech certificate for a map
attributed there to Fable; lines 31--38 expressly say that an original
model-session transcript was not recovered there.  The JSON supplied with
this audit is a new exact replay record, not that missing discovery transcript.

\subsection{The original polynomial, its determinant, and its collision}
Retain the original ordered polynomial coordinates $(x,y,t)$ and define
$F:\mathbb C^3\longrightarrow\mathbb C^3$ by
\begin{align}
 F_1&=(1+xy)^3t+y^2(1+xy)(4+3xy),\label{eq:s6b:F1}\\
 F_2&=y+3x(1+xy)^2t+3xy^2(4+3xy),\label{eq:s6b:F2}\\
 F_3&=2x-3x^2y-x^3t.\label{eq:s6b:F3}
\end{align}
The symbol $t$ in these three equations is a spatial coordinate, never fluid
time.  Fluid time below is $s$.  To compare with a convention using spatial
coordinates $(z_1,z_2,z_3)$, use the explicitly defined linear isometry
$R(z_1,z_2,z_3)=(x,y,t)=(z_1,z_2,z_3)$.  Its derivative is $I_3$,
determinant is $1$, and its pullback sends
$(\partial_x,\partial_y,\partial_t)$ to
$(\partial_{z_1},\partial_{z_2},\partial_{z_3})$ in the same order.
Thus $F\circ R$ is an exact renaming preserving every coefficient,
derivative, volume form and Euclidean operator.  We continue to use $(x,y,t)$.

Here is a determinant proof retaining also the locus $x=0$.  On $x\ne0$,
put $v=1+xy$ and $h=2-3xy-x^2t$.  Direct substitution gives
\[
 F_1=\frac{v^2+v-v^3h}{x^2},\qquad
 F_2=\frac{4v+2-3v^2h}{x},\qquad F_3=xh.
\]
The coordinate-change determinant is
\[
 \det\frac{\partial(x,v,h)}{\partial(x,y,t)}
 =\det\begin{pmatrix}1&0&0\\y&x&0\\-3y-2xt&-3x&-x^2\end{pmatrix}
 =-x^3.
\]
Writing $A=v^2+v-v^3h$ and $B=4v+2-3v^2h$, the other determinant is
\[
 \det\frac{\partial(F_1,F_2,F_3)}{\partial(x,v,h)}
 =x^{-3}\det\begin{pmatrix}
 -2A&2v+1-3v^2h&-v^3\\
 -B&4-6vh&-3v^2\\ h&0&1
 \end{pmatrix}=\frac2{x^3}.
\]
Indeed the displayed numerator expands to
\[
 (4-6vh)(-2v^2-2v+3v^3h)
 +(2v+1-3v^2h)(4v+2-6v^2h)=2.
\]
Consequently $\det DF=-2$ on $x\ne0$.  Both sides are polynomials in the
original coordinates, so equality on this open set is a polynomial identity
and includes $x=0$.  Independently, substitution at $x=0$ into the original
Jacobian gives determinant $-2$ there as well.

The three retained points
\[
 p_0=(0,0,-1/4),\quad p_+=(1,-3/2,13/2),\quad
 p_-=(-1,3/2,13/2)
\]
are distinct and have common image $(-1/4,0,0)$.  At $p_0$ this follows
directly from the original equations.  At $p_\pm$, one has
$xy=-3/2$, $1+xy=-1/2$, $4+3xy=-1/2$, $y^2=9/4$ and $t=13/2$.
Thus $F_1=-13/16+9/16=-1/4$.  At $p_+$,
$F_2=-12/8+39/8-27/8=0$ and $F_3=2+9/2-13/2=0$.
At $p_-$ each summand in $F_2,F_3$ changes sign, so both remain zero.
This is an exact algebraic certificate about the displayed map, independent
of authorship or discovery claims.
In fact these are the entire fibre, including over $\mathbb C$.
For a point with image $(-1/4,0,0)$ and $x=0$, the original second equation
gives $y=0$ and the first then gives $t=-1/4$.
For $x\ne0$, $F_3=xh=0$ gives $h=0$; the retained equations give
$xF_2=4v+2=0$, hence $v=-1/2$, and
$x^2F_1=v^2+v=-1/4$.  Since $F_1=-1/4$, this forces $x^2=1$.
Then $y=(v-1)/x=-3/(2x)$ and
$t=(5-3v)/x^2=13/2$.  Both possibilities are exactly $p_+$ and $p_-$.

\subsection{A polynomial incompressible velocity and an escaping trajectory}
Restrict $F$ to $\mathbb R^3$ and let $J=DF$ in the original coordinate order.
Since $\det J=-2$, the field $W=J^{-1}e_1$ is globally polynomial.  Its full
components are
\begin{align}
 W_x={}&3x^6yt+9x^5y^2+3x^5t+3x^4y-4x^3,\\
 W_y={}&3x^4yt+9x^3y^2+3x^3t+3x^2y-3x,\\
 W_t={}&-9x^5yt^2-45x^4y^2t-9x^4t^2-54x^3y^3
          -30x^3yt-27x^2y^2+9x^2t+21xy+1.
 \label{eq:s6b:W}
\end{align}
These components can be verified without an inverse-matrix convention:
cofactor expansion gives
\begin{equation}
 W=\frac{\nabla F_2\times\nabla F_3}{\det J}
   =-\frac12\nabla F_2\times\nabla F_3
   =\operatorname{curl}\!\left(-\frac12 F_2\nabla F_3\right).
 \label{eq:s6b:potential}
\end{equation}
The dot product of $\nabla F_1$ with the numerator is $\det J$; the dot
products with $\nabla F_2,\nabla F_3$ are zero.  Hence $JW=e_1$ exactly,
proving the inverse claim.  Expanding the three cross-product components
gives precisely \eqref{eq:s6b:W}.  The identity
$\operatorname{curl}(a\nabla b)=\nabla a\times\nabla b$ follows by
the product rule and cancellation of the symmetric second derivatives of
$b$.  Mixed derivatives likewise give $\operatorname{div}\operatorname{curl}=0$.
In particular $\operatorname{div}W=0$.  An independent component check is
\begin{align*}
 \partial_xW_x&=18x^5yt+45x^4y^2+15x^4t+12x^3y-12x^2,\\
 \partial_yW_y&=3x^4t+18x^3y+3x^2,\\
 \partial_tW_t&=-18x^5yt-45x^4y^2-18x^4t-30x^3y+9x^2,
\end{align*}
whose sum is zero coefficient by coefficient.

For $s<1/4$ use the positive real square root and define
\begin{equation}
 \gamma(s)=\bigl((1-4s)^{-1/2},
          -\tfrac32(1-4s)^{1/2},\tfrac{13}{2}(1-4s)\bigr).
 \label{eq:s6b:trajectory}
\end{equation}
The entire displayed interval is part of its definition; in particular
$\gamma(0)=p_+$.  Along the curve, write $X=(1-4s)^{-1/2}$.
Then $y=-3/(2X)$ and $t=13/(2X^2)$.  Substituting into
\eqref{eq:s6b:W} gives
$W(\gamma(s))=(2X^3,3X,-26)$, which is exactly $\gamma'(s)$.
The restricted original map gives
\[
 F(\gamma(s))=(-1/(4X^2),0,0)=(s-1/4,0,0).
\]
Equivalently, this last equality and $J\gamma'=e_1$ give the same trajectory
identity because $J$ is invertible everywhere.  As $s\uparrow1/4$, the first
coordinate tends to $+\infty$.  Therefore this integral curve has no
continuous extension to time $1/4$ in $\mathbb R^3$: any such extension
would have a finite first coordinate.  This is finite-time escape of a
particle trajectory in a stationary smooth polynomial vector field.

\subsection{Exact Euclidean Navier--Stokes residual}
Use the Euclidean Laplacian
$\Delta=\partial_x^2+\partial_y^2+\partial_t^2$ and viscosity $\nu>0$.
A stationary unforced Navier--Stokes solution with velocity $W$ would satisfy
\begin{equation}
 (W\cdot\nabla)W-\nu\Delta W=-\nabla p.
 \label{eq:s6b:steady-ns}
\end{equation}
The right-hand side has zero curl, even for a distributional pressure,
because distributional mixed derivatives commute.  Direct differentiation
of the retained polynomial gives
\begin{equation}
 \left.\operatorname{curl}\bigl((W\cdot\nabla)W-\nu\Delta W\bigr)
 \right|_{(x,y,t)=(0,0,t)}=(0,0,-18\nu t).
 \label{eq:s6b:curl-obstruction}
\end{equation}
Here is a local calculation of that identity in addition to the exact replay.
Along this axis $W=(0,0,1)$ and the only possibly nonzero entry of $DW$ is
$\partial_xW_y=-3$; $\partial_xW_t=21y$ vanishes there.  Thus
$(W\cdot\nabla)W=\partial_tW=0$ there.  Differentiating this advection
expression before restriction gives zero first derivatives there: the
product-rule term from $\partial_xW_y=-3$ multiplies $\partial_yW=0$,
and the remaining term $\partial_{xt}W$ also vanishes; derivatives in
$y,t$ vanish by the same displayed polynomial.
Consequently its curl is zero there.  The viscous curl is read from
\begin{align*}
 W_x&=-4x^3+3x^4y+3x^5t+9x^5y^2+3x^6yt,\\
 W_y&=-3x+3x^2y+3x^3(t+3y^2)+3x^4yt,\\
 W_t&=1+21xy+x^2(9t-27y^2)+x^3(-30yt-54y^3)\\
    &\hspace{14mm}+x^4(-9t^2-45y^2t)-9x^5yt^2.
\end{align*}
At $x=y=0$,
$\partial_y\Delta W_t-\partial_t\Delta W_y=0$,
$\partial_t\Delta W_x-\partial_x\Delta W_t=0$, and
$\partial_x\Delta W_y-\partial_y\Delta W_x=18t$.
These are respectively the three components of $\operatorname{curl}\Delta W$.
They prove \eqref{eq:s6b:curl-obstruction}, including its sign.
At the single point $(0,0,1)$ the residual curl is $(0,0,-18\nu)\ne0$.
Thus no pressure makes $W$ an unforced stationary Euclidean Navier--Stokes
solution for any positive viscosity.

There is also an independent initial-data obstruction.  $W$ has infinite
Euclidean kinetic energy.  For completeness, a nonzero polynomial $P$ of
degree $d$ cannot lie in $L^2(\mathbb R^3)$: its top homogeneous part is
nonzero on some open spherical patch, and is bounded in absolute value
below by a positive constant on a smaller compact patch.  On the associated
cone, for all sufficiently large radii, lower-degree terms are bounded by
half the leading term, so $|P(r\theta)|\ge c r^d$.
The polar-coordinate integral then dominates
$c^2\int_R^\infty r^{2d+2}\,dr=\infty$ times a positive solid angle.
Apply this to a nonzero component of $W$; its first component has degree $8$.
Moreover $W$ is not unit-periodic: on $y=t=0$ its first component is
$-4x^3$, whose values at $x=0,1$ differ.  None of these facts establishes
anything about every possible velocity constructed from $F$; they establish
the precise limitations of the displayed $W$.

\subsection{Exact local-isometry correspondence and incomplete metric}
There is a stronger geometric bridge.  Equip the same underlying
$\mathbb R^3$ with
\begin{equation}
 g=J^{\mathsf T}J,
 \qquad g_{ab}=\sum_{i=1}^3(\partial_aF_i)(\partial_bF_i).
 \label{eq:s6b:metric}
\end{equation}
This is positive definite because $J$ is invertible.  Its determinant is $4$,
so its positive volume density is $2\,dx\,dy\,dt$; the signed pullback of
the target orientation is instead $F^*(dX_1\wedge dX_2\wedge dX_3)
=-2\,dx\wedge dy\wedge dt$.  The real inverse function theorem gives
local coordinates $X_i=F_i$ around every point.  In exactly these coordinates
the metric is $dX_1^2+dX_2^2+dX_3^2$.  Thus $F$ is a flat local isometry,
with the specified orientation reversal; it is not asserted to be a global
coordinate chart.
The metric change already appears at the original origin:
\[
 J(0)=\begin{pmatrix}0&0&1\\0&1&0\\2&0&0\end{pmatrix},
 \qquad g(0)=\operatorname{diag}(4,1,1).
\]
For a direct scalar diffusion comparison, the target coordinate $X_1$
is Euclidean harmonic, while $F^*X_1=F_1$ has
$\Delta F_1(0)=8$, from its exact quadratic term $4y^2$.
The metric correspondence therefore does not intertwine the two
ordinary Euclidean Laplacians.

Let $\Omega=F(\mathbb R^3)$, an open set, and let a smooth target velocity
$v(s,X)$ and pressure $p(s,X)$ be defined on an open time interval times
$\Omega$.  Define
\begin{equation}
 u(s,x,y,t)=J(x,y,t)^{-1}v(s,F(x,y,t)),\qquad
 p_g(s,x,y,t)=p(s,F(x,y,t)).
 \label{eq:s6b:pullback}
\end{equation}
This is a globally defined pullback map on those fields.  In each of the
local coordinates $X=F(x,y,t)$ its components and metric are precisely
$v$ and the Euclidean metric.  Consequently the following operator
identities hold on every local chart, and hence globally:
\begin{align*}
 \operatorname{div}_g u&=(\operatorname{div}v)\circ F,\\
 \nabla^g_u u&=J^{-1}((v\cdot\nabla_X)v)\circ F,\\
 \operatorname{grad}_g p_g&=J^{-1}(\nabla_Xp)\circ F,\\
 (\Delta_{H,g}u^\flat)^\sharp
       &=J^{-1}(-\Delta_Xv)\circ F.
\end{align*}
Here $\Delta_{H,g}=d\delta_g+\delta_gd$ is the nonnegative Hodge Laplacian;
on Euclidean one-forms it is the negative componentwise Laplacian.
One may verify the identities directly in $X$ coordinates: the metric
components are constant, Christoffel symbols vanish, the divergence is
$\sum_i\partial_{X_i}v_i$, and the Hodge Laplacian is
$-\sum_i\partial_{X_i}^2$ on each component.  Coordinate tensor
transformation gives exactly the displayed $J^{-1}$ factors.
The $g$-divergence also equals ordinary coordinate divergence because its
volume density is the constant $2$:
$\operatorname{div}_gu=\frac12\sum_a\partial_a(2u^a)
=\operatorname{div}u$.

It follows by substitution, retaining the same $s$ and $\nu$, that pullback
carries the Euclidean equation
$\partial_sv+(v\cdot\nabla_X)v=-\nabla_Xp+\nu\Delta_Xv$
to the exact intrinsic equation
\[
 \partial_su+\nabla^g_u u
 =-\operatorname{grad}_gp_g-\nu(\Delta_{H,g}u^\flat)^\sharp,
 \qquad \operatorname{div}_gu=0.
\]
This is a proved morphism between specified PDE solution spaces.
Conversely a smooth field $u$ descends under this morphism precisely when
$J(x)u(x)$ is constant on each fibre of $F$; a pressure descends precisely
when it too is constant on each fibre.  Necessity follows from
\eqref{eq:s6b:pullback}.  For sufficiency these equalities define a field
and pressure on $\Omega$; local inverses of $F$ prove smoothness, and
fibre constancy makes the local definitions agree.  The local operator
identities then prove the equations for the descended field.  There is
therefore a full local correspondence and a stated exact global descent
condition, not an assertion of a nonexistent global inverse.

In particular the target solution $v=e_1$, $p=0$ pulls back to $u=W$.
It is parallel for $g$, has $|W|_g=1$, and solves this metric equation for
every $s$ and every $\nu>0$.  Its metric energy is also infinite:
$\int_{\mathbb R^3}|W|_g^2\,d\operatorname{vol}_g
=2\int_{\mathbb R^3}1\,dx\,dy\,dt=\infty$.
Nevertheless $g$ is incomplete.  The
$g$-speed of \eqref{eq:s6b:trajectory} is
$|J\gamma'|=|e_1|=1$.  For any sequence $s_n\uparrow1/4$, the tail of
the curve has length $|s_n-s_m|$, so the points $\gamma(s_n)$ are Cauchy
in the $g$ distance.  They cannot converge to a point of $\mathbb R^3$:
the first coordinate is unbounded, and a smooth positive-definite metric
induces the ordinary local topology.  To check the latter assertion at a
putative limit, take a small closed Euclidean ball around it.  The minimum
eigenvalue of $g$ on this compact ball is positive; every curve reaching
the ball's boundary has $g$ length at least that lower length factor
times the Euclidean radius.  Thus sufficiently small metric balls lie
in any chosen Euclidean neighborhood of the limit.  This contradicts
the unbounded coordinate of the sequence.
The metric field remains smooth at every finite spatial point and at every
time; particle escape in this incomplete metric is not finite-time loss
of smoothness in the complete Euclidean spatial geometry.

\subsection{A compactly supported incompressible initial-data map}
An explicit localization retains a genuine map to smooth admissible spatial
data.  Define $b(r)=e^{-1/r}$ for $r>0$ and $b(r)=0$ for $r\le0$, and
\[
 \eta(r)=\frac{b(4-r)}{b(4-r)+b(r-1)},\qquad
 \chi(z)=\eta(|z-p_+|^2/R^2),\quad R>0.
\]
The denominator is positive for every real $r$.  Each derivative of $b$
on the positive side is an exponential times a polynomial in $1/r$,
which tends to zero as $r\downarrow0$; hence $b$ and $\eta$ are smooth.
The function $\chi$ equals $1$ on the closed radius-$R$ ball about $p_+$
and $0$ outside the radius-$2R$ ball.  Set
\begin{equation}
 U_\chi=-\frac12\operatorname{curl}(\chi F_2\nabla F_3)
 =\chi W-\frac12 F_2(\nabla\chi\times\nabla F_3).
 \label{eq:s6b:cutoff}
\end{equation}
The product rule proves the second equality.  All factors are smooth,
and the potential has compact support, so $U_\chi$ is smooth and compactly
supported.  Its divergence is zero by commuting mixed derivatives.
On the plateau $\chi\equiv1$ its derivative is zero, giving
$U_\chi=W$ exactly.  Every derivative of $U_\chi$ is bounded and
compactly supported; in particular the field is in every Sobolev space
of nonnegative integral order, has finite energy, and decays faster than
every inverse power.  It is thus an explicit map to smooth finite-energy
divergence-free initial data for the Euclidean problem.

The autonomous particle equation for this compact field is complete.
Indeed $|U_\chi|\le M$ globally and its first derivatives are bounded,
so it is globally Lipschitz.  On every finite time interval a trajectory
stays within distance $M$ times the length of that interval from its
initial point; standard local ODE existence and uniqueness extend it
past any finite endpoint in that compact ball.  This completeness
statement concerns its autonomous ODE and makes no assertion about
the subsequent nonlinear Navier--Stokes evolution of $U_\chi$.
The escaping curve $\gamma$ leaves the plateau in finite time before
$1/4$ because its first coordinate diverges.  Thus the exact local
agreement does not preserve the escaping final segment.

\subsection{Retained cusp Fourier data mapped to exact periodic solutions}
The source
\path{supporting_materials/workbench/research/cusp_spectral_geometry.tex},
lines 83--185, specifies an additional Euclidean metric on a regular
four-dimensional real torus.  Retain all original coefficients:
\[
 P_T=\begin{pmatrix}
 6\alpha_T&\eta_T&1&0\\6m_T&L_T&0&0\\
 c_T&\alpha_T&0&1\\-q_T&m_T&0&0
 \end{pmatrix},\qquad D_T=L_Tq_T+6m_T^2>0,
 \quad F_T=\mathbb R^4/P_T\mathbb Z^4.
\]
Its oriented determinant is $-D_T$.  For the original dual label
$k=(r,s_0,n_1,n_2)\in\mathbb Z^4$ define
\[
 R_T=r-6\alpha_Tn_1-c_Tn_2,\qquad
 S_T=s_0-\eta_Tn_1-\alpha_Tn_2.
\]
Here $s_0$ denotes the source's second integer Fourier label, solely to
avoid collision with fluid time $s$; the map $s_{\rm source}\mapsto s_0$
is the identity on that integer, with no change to the label or lattice.
Solving $P_T^{\mathsf T}\ell=k$ gives exactly
\[
 \ell=\left(n_1,\frac{m_TR_T+q_TS_T}{D_T},n_2,
                  \frac{-L_TR_T+6m_TS_T}{D_T}\right).
\]
The first and third components come from the last two equations; the other
two solve
$\begin{pmatrix}6m_T&-q_T\\L_T&m_T\end{pmatrix}
(\ell_2,\ell_4)^{\mathsf T}=(R_T,S_T)^{\mathsf T}$,
whose determinant is $D_T$.  This proves the inverse-transpose formula,
including every sign.  The character
$\exp(2\pi i\langle\ell,x\rangle)$ is well-defined on $F_T$ since
$\langle\ell,P_Tn\rangle=k\cdot n\in\mathbb Z$.
The nonnegative scalar Laplacian $\Delta_T$ therefore has the eigenvalue
$4\pi^2|\ell|^2$ on this character by two differentiations.

Fix any nonzero original label $k$ and put $a=|\ell|>0$.
Consider the exact closed sector of functions
$h([x])=f(\langle\ell,x\rangle\bmod1)$ with smooth real $1$-periodic $f$.
The phase map is surjective onto the circle: its lift is a nonzero real
linear functional, so takes every real value.  Hence this representation
determines $f$ uniquely.  On this sector define
\[
 \mathcal I_k h(z_1,z_2,z_3)=e_2 f(z_1),
 \qquad (z_1,z_2,z_3)\in\mathbb R^3/\mathbb Z^3.
\]
This is an injective linear map with fully specified source and target.
If $\Delta_3=-\sum_{j=1}^3\partial_{z_j}^2$ is the nonnegative
componentwise operator on the target, the chain rule proves
\[
 \Delta_T h=-a^2f''(\langle\ell,x\rangle),\qquad
 \Delta_3\mathcal I_kh=e_2(-f''(z_1))
   =a^{-2}\mathcal I_k\Delta_T h.
\]
Thus the source heat solution $h(\tau)=e^{-\nu\tau\Delta_T}h_0$
gives the target velocity
\[
 u(s,z)=\mathcal I_k h(s/a^2)(z),\qquad p(s,z)=0,
 \quad s\ge0.
\]
It satisfies $\partial_su=-\nu\Delta_3u$.  Also
$\operatorname{div}u=\partial_{z_2}f(z_1)=0$ and
$(u\cdot\nabla)u=f(z_1)\partial_{z_2}(e_2f(z_1))=0$.
Therefore it solves the full unforced periodic three-dimensional
Navier--Stokes equation with viscosity $\nu$, exactly.

For an explicit global regularity proof in this sector, write
$f_0(\theta)=\sum_{n\in\mathbb Z}c_ne^{2\pi in\theta}$ with
$c_{-n}=\overline{c_n}$.  Integration by parts $N$ times in the Fourier
coefficient integral bounds $|c_n|\le C_N|n|^{-N}$ for every $N$ and
$n\ne0$.  The target solution is
\[
 u(s,z)=e_2\sum_{n\in\mathbb Z}
        c_n e^{-4\pi^2\nu n^2s}e^{2\pi in z_1}.
\]
After any prescribed finite number of time and space derivatives, choose
$N$ larger than the resulting power of $n$ plus $2$.  The differentiated
series then converges absolutely and uniformly for every $s\ge0$.
It follows that this solution is smooth for all nonnegative time; its
energy on the unit torus is finite, and its mean is zero when $c_0=0$.
No singularity is obtained by this exact Fourier-sector morphism.

The source gauge operator at
\path{supporting_materials/workbench/research/cusp_gauge_operators.tex},
lines 64--95, is the Hessian $g_{\mathrm{YM}}^{-2}d^*d$ at a product
connection.  On a co-closed one-form, $d^*d$ agrees with the Hodge
Laplacian because $dd^*$ vanishes.  This identifies its differential
expression with the positive linear diffusion operator; it supplies no
nonzero convection term.  In the exact periodic sector above that
term was computed explicitly and is zero.

\subsection{Reproducibility and precise conclusion}
The exact rational-polynomial script
\path{scripts/s6_bridge_check.py} retains the original $F$,
checks $\det DF=-2$, all three collision images, $JW=e_1$, the vector
potential, divergence, trajectory, residual curl, the full retained cusp
matrix determinant and inverse transpose.  It writes
\path{research/s6_bridge_audit.json}, including SHA-256
hashes of the seven source files used.  All assertions passed with
SymPy 1.13.1 on 8 September 2026; a separate agent independently replayed
the velocity, trajectory, potential and residual-curl calculations and
confirmed their signs.
The fixed source hashes are packaged in
\path{sources/s6_source_hashes.json}.
Default replay uses that recorded provenance and requires no sibling source
repository; it does not claim to reread the seven source files.
The optional \texttt{--source-root} argument additionally rereads those
original files and fails on a hash mismatch.  The result records which
of these two provenance modes was used.

The exact relations established here are a polynomial developing map,
its divergence-free inverse-Jacobian field, its flat local-isometry PDE
pullback, its compactly supported divergence-free localization, and its
cusp Fourier-sector map to classical periodic Navier--Stokes solutions.
The escaping polynomial field fails the unforced Euclidean PDE by the
nonzero curl in \eqref{eq:s6b:curl-obstruction}; the metric pullback is
smooth but has incomplete spatial geometry; the localized fields are
admissible initial data without an established singular evolution; and
the Fourier-sector solutions are globally smooth.  Accordingly these
calculations provide no verified disproof of the classical
three-dimensional Navier--Stokes existence-and-smoothness assertion.
They also do not assert absence of other mathematical relations or rule
out other constructions from the retained objects.
